\documentclass[journal]{IEEEtran}
\usepackage{graphicx} % Required for inserting images
\usepackage{amsmath} % For \text in math mode
\usepackage{float}
\usepackage{hyperref}
\usepackage{algorithm}
\usepackage{algpseudocode}
\usepackage{amssymb}
\usepackage{booktabs}
\usepackage{siunitx}
\usepackage{tikz}
\usepackage{adjustbox}
\usepackage{multirow}


\usepackage{cleveref}
\crefname{section}{Section}{Sections}
% \crefname{subsection}{Subsection}{Subsections}

\usetikzlibrary{arrows.meta,positioning}

\title{HIT-ASR: Hierarchical Transformer Routing for Adaptive ASR Expert Selection}

\author{A.~Samil~Namli$^{*}$, 
        Huseyin~Karaca$^{*}$,
        and Suleyman S.~Kozat,~\IEEEmembership{Senior~Member,~IEEE}
\thanks{$^{*}$A. S. Namli and H. Karaca contributed equally to this work.}
\thanks{A. S. Namli, H. Karaca, and S. S. Kozat are with Bilkent University, Ankara, Turkiye. A. S. Namli and H. Karaca are also with Turk Telekom, Ankara, Turkiye (e-mails: \{samil.namli, huseyin.karaca\}@bilkent.edu.tr; kozat@ee.bilkent.edu.tr).}%
\thanks{This work is supported by The Scientific and Technological Research Council of Turkiye (TUBITAK) 1515 Frontier R\&D Laboratories Support Program for Turk Telekom 6G R\&D Lab under project number 5249902, and by the TUBITAK 2210-A program.}%
}
\date{}

\begin{document}

\maketitle 

\begin{abstract}
Adaptive expert selection is essential for real-world Automatic Speech Recognition (ASR) because no single model performs optimally across all acoustic environments and speaker variations. Conventional routing methods often collapse utterance features into a single sequence-level embedding, inadvertently obscuring the fine-grained temporal dynamics critical for identifying the most suitable expert. In this paper, we introduce HIT-ASR, a hierarchical transformer routing framework that operates directly on the frame-level encoder representations of multiple pretrained ASR systems. By preserving intra-clip temporal structure, HIT-ASR can dynamically select the best expert per clip. The router leverages a two-stage transformer architecture with a cross-attention bridge to model within-expert temporal evidence and cross-expert relationships. Furthermore, we propose a Word Error Rate (WER)-aware training objective infused with a soft cross-entropy penalty to explicitly optimize transcription performance. Extensive evaluations on a synthetic regime-switch dataset and the real-world AMI and VoxPopuli corpora demonstrate that HIT-ASR consistently outperforms static single-model approaches and pooling-based selection baselines, significantly reducing WER. By avoiding the computational bottleneck of autoregressive decoding for all candidate models and executing only the selected expert's decoder, HIT-ASR maintains an inference cost comparable to a single ASR model.
\end{abstract}

\begin{IEEEkeywords}
Speech recognition, automatic speech recognition (ASR),
speech-to-text, model selection, adaptive routing, mixture of
experts, transformer, frame-level embeddings, ensemble of pretrained
models.
\end{IEEEkeywords}


\section{Introduction}
\label{sec:introduction}

\subsection{Background}

Automatic speech recognition (ASR, also referred to as speech-to-text) has progressed significantly with deep sequence models \cite{hinton_deepn,dahl_context,deng_conf}. The Transformer architecture \cite{vaswani2017attention} and its convolution-augmented variant, the Conformer \cite{gulati2020conformer}, now achieve state-of-the-art accuracy on benchmarks such as LibriSpeech, outperforming recurrent and CNN models \cite{deng_conf,bigssl}. Self-supervised pretraining \cite{baevski2020wav2vec,hsu2021hubert,chen2022wavlm} further boosts accuracy, and streaming architectures like the RNN-Transducer \cite{Graves2012SequenceTW}, and its Transformer-based extensions enable low-latency transcription \cite{gong_advanced,lee_mle}. However, these systems are static because they use a fixed model for all inputs. In realistic deployments, acoustic conditions such as noise, channel and speaker can vary unpredictably over time. A static model may be suboptimal across domains or may waste capacity on easy inputs. Consequently, no single ASR system is optimal across all scenarios \cite{li2014overviewrobust,cornell23_chime,cornell24_chime}. Different models perform better in different acoustic regimes, forming a complementary set of experts. Using this diversity requires an adaptive mechanism to select the most suitable model for each input, rather than committing to a fixed system.

A key difficulty in this setting is predicting the performance of candidate models without explicitly executing them \cite{kerschke2019algorithm_selection}. Observable features of the audio signal, such as acoustic statistics, temporal dynamics, and learned representations, provide useful cues for this purpose \cite{baevski2020wav2vec,hsu2021hubert,li2014overviewrobust,bell2021adaptation}. Several adjacent ASR confidence and selection approaches compress an utterance into a sequence-level score or embedding before making a downstream decision \cite{kumar2020utterance_confidence,swarup2019acoustic_hypothesis_embeddings,li2021rebm}. Such pooling can obscure temporal structure that may be relevant for routing, because speech acoustics and interaction conditions evolve over time \cite{li2014overviewrobust,li2023asr_emotional_speech}. Within a single clip, the speaker, the channel, the background, and the spoken content are rarely homogeneous \cite{li2014overviewrobust,bell2021adaptation,cornell23_chime,vinnikov2024notsofar1}. For example, a single meeting utterance may contain both overlapping speech and clean segments \cite{yoshioka2018overlapped_meetings,cornell23_chime,cornell2024chime8,vinnikov2024notsofar1}. A long-form public or meeting recording may alternate between near-field, distant, and reverberant segments, while speaking rate or emotion may also change across the conversation. The expert that transcribes such a clip best may depend strongly on these intra-clip dynamics, rather than only on a single global summary \cite{kumar2020utterance_confidence,li2023asr_emotional_speech}.

Ignoring such intra-clip variations can  cause suboptimal routing decisions, especially in long-form audio. This motivates the development of adaptive ASR frameworks that explicitly model the frame-level temporal structure of each clip when routing among experts. Here, we introduce HIT-ASR, a hierarchical transformer router that operates directly on the frame-level encoder representations of each clip, rather than on pooled features. In essence, the system first processes a long audio stream into distinct clips. The router then uses frame-level embeddings extracted from candidate ASR encoders to select the most suitable ASR expert for each clip. A Transformer-based routing function models the temporal dynamics within these frames to produce expert selection scores. At inference time, the system decodes each clip using only the selected ASR expert, so latency remains comparable to that of a single model per clip. By avoiding the need to execute the computationally expensive autoregressive decoding steps for all candidate models, HIT-ASR circumvents the primary bottleneck of ensemble methods.

 
\subsection{Related Work}
\label{sec:related_work}

Modern end-to-end ASR often uses attention-based encoder-decoder or transducer models \cite{deng_conf}. The original Transformer \cite{vaswani2017attention} eliminated recurrence in sequence-to-sequence tasks and has been widely adapted for speech. In ASR, the Conformer architecture \cite{gulati2020conformer} augments self-attention with convolutional modules to better model local speech structure. It achieved near-state-of-the-art LibriSpeech WERs (2.1\% clean, 4.3\% other) while remaining compact \cite{gulati2020conformer}. Self-supervised pretraining methods, such as wav2vec 2.0 \cite{baevski2020wav2vec}, learn useful speech representations from unlabeled audio, enabling strong ASR performance even with limited labeled data. Streaming extensions include the RNN-Transducer \cite{Graves2012SequenceTW}, Transformer-Transducer models that incorporate attention-based components into the transducer framework, and chunk-based Conformers designed for real-time recognition. These models primarily focus on improving accuracy and, in streaming settings, reducing latency, but they remain static because all inputs are processed by the same network. Our work is orthogonal to these approaches. Rather than modifying the internal architecture of a single ASR model, we route each audio clip to one of several pretrained ASR systems and decode only with the selected model. As a result, the decoding cost remains comparable to that of a single ASR model, while preserving the diversity of the expert pool.

Adaptive expert selection is also closely related to the broader literature on mixture-of-experts (MoE) \cite{jacobs1991,jordan1994hierarchial,yuksel2012twenty}. Modern sparse deep MoE models enabled enormous capacity by routing each token to a sparse subset of experts \cite{shazeer2017outrageously}. This idea was scaled in GShard \cite{gshard} and Switch Transformer \cite{fedus2021switch}, where models with hundreds of billions of parameters were trained with sparsely-gated MoE layers at constant cost. Beyond NLP, MoE has been applied in vision and other domains as well \cite{yuksel2012twenty,riquelme2021vmoe}. There is growing work on adaptive or multi-domain ASR. SpeechMoE \cite{you2021speechmoe} adds MoE layers to a CTC model with a learned embedding for the router. This approach achieved substantial CER reductions (7--23\% relative) over a static model by dynamically routing frames to experts \cite{you2021speechmoe}. A follow-up SpeechMoE2 \cite{speechmoe2} improved routing by feeding domain- and accent-embeddings to the router. \cite{hu2023moeconformer} introduced a Switch-Conformer MoE, a multilingual
streaming Conformer with MoE FFN layers, reporting an $11.9\%$ relative WER drop while activating only $53\%$ of parameters at inference. Our router differs from these methods in both granularity and scope. Whereas SpeechMoE-style layers route individual frames within a single trained model, our approach selects among complete pretrained ASR systems for an entire audio clip. The routing decision is based on frame-level encoder outputs extracted from those systems.

More broadly, conditional computation and dynamic neural networks adapt inference to the input by selectively activating parts of the computation graph or terminating inference early \cite{han2022dynamic,scardapane2024conditional,laskaridis2021adaptive}. Representative mechanisms include dynamic skipping of recurrent updates or residual blocks \cite{campos2018skip,wu2018blockdrop,wang2018skipnet}, early-exit architectures with intermediate classifiers \cite{teerapittayanon2016branchynet,huang2018msdnet,xin2020deebert}, and adaptive-depth sequence models that vary computation across inputs or positions \cite{graves2016adaptive,dehghani2019universal,elbayad2020depth}. For long-sequence transformers, Routing Transformer learns content-based sparse attention patterns with sub-quadratic complexity \cite{roy-efficient}. Sparse expert layers such as BASE Layers similarly allocate tokens to experts within a single model \cite{lewis2021base}. By contrast, our setting routes among independently pretrained ASR systems with no shared parameters, so the router learns inter-model selection rather than intra-network conditional computation. 


% Adaptive expert selection is also closely related to the broader
% literature on mixture-of-experts (MoE) and conditional computation.
% In classical MoE models, a gating network assigns input-dependent
% weights to multiple experts, and both the experts and the gating
% mechanism are trained jointly within a unified architecture.
% Recent developments in sparse and large-scale MoE systems further
% improve efficiency by activating only a subset of parameters for each
% input. However, these approaches typically assume that experts share
% a common backbone and are optimized end-to-end under a unified
% objective. In contrast, our setting considers independently pretrained
% ASR systems as fixed experts, with no parameter sharing or joint
% training. The routing mechanism operates as a separate model that
% learns to predict expert suitability based on input characteristics,
% making it possible to incorporate heterogeneous models without
% modifying their internal architectures.

% A related direction explores confidence estimation and quality
% prediction for ASR outputs. Confidence scores derived from posterior
% probabilities or auxiliary prediction models have been used for
% hypothesis selection, rescoring, and fallback strategies. Similarly,
% domain classification and environment detection methods attempt to
% route inputs to specialized models based on coarse categorical
% attributes such as noise type or speech domain. While effective in
% certain scenarios, these approaches either require partial decoder
% execution to obtain confidence measures or rely on predefined domain
% boundaries that may not capture fine-grained variability in acoustic
% conditions. In contrast, feature-based selection methods aim to
% predict expert performance directly from input representations,
% enabling lightweight routing decisions without running multiple ASR
% systems.

% Most existing adaptive selection approaches formulate the problem as
% an instance-wise decision task, where each audio segment is processed
% independently. This assumption neglects the temporal structure of
% speech signals, where consecutive segments are often correlated and
% share common acoustic characteristics. In long-form recordings, such
% as conversations and meetings, factors such as speaker identity,
% background noise, and channel conditions tend to persist over time or
% change gradually. Ignoring these temporal dependencies can limit the
% effectiveness of selection mechanisms, leading to unstable or
% suboptimal routing decisions. In this work, we address this limitation
% by modeling adaptive ASR selection as a sequential decision problem
% and incorporating temporal context into the routing process, enabling
% more consistent and context-aware expert selection across audio
% streams.


\subsection{Contributions}

Our contributions are summarized as follows:

\begin{itemize}

\item We introduce HIT-ASR, an adaptive ASR expert-selection framework that routes each audio clip to the most suitable pretrained ASR system using frame-level encoder representations. Unlike pooling-based routers, the proposed framework preserves intra-clip temporal structure for routing decisions.

\item We formulate adaptive ASR selection as a per-clip classification problem over a fixed pool of pretrained ASR experts and propose a hierarchical transformer router that models both within-expert temporal evidence and cross-expert relationships.

\item We develop a WER-aware training objective that directly optimizes the expected transcription error while using auxiliary hard-label and soft-label supervision to improve training stability. The framework can also support alternative transcription-quality metrics without changing the routing architecture.

\item We design the framework for efficient inference: the router uses encoder hidden states associated with the candidate ASR systems, while only the selected expert is fully decoded for each input. As a result, the decoding cost remains comparable to running a single ASR model.

\item We present a routing-model-agnostic formulation in which the hierarchical transformer can be replaced by alternative sequence learners, such as recurrent or convolutional models, without changing the overall adaptive ASR pipeline.

\item We share all our codes for further research and reproducibility at \url{https://github.com/huseyin-karaca/hit-asr}.

\end{itemize}


\subsection{Organization}

The remainder of the paper is organized as follows.
\cref{sec:problem} formalizes the per-clip adaptive ASR selection
problem. \cref{sec:workdone} presents the proposed framework,
including the frame-level feature representation, the hierarchical
transformer router, and the training objective.
\cref{sec:experiments} describes the experimental setup and empirical
evaluation. Finally, \cref{sec:conclusion} concludes the paper.

\section{Problem Description}
\label{sec:problem}

We consider the problem of adaptive transcription at the level of individual audio clips. Let $\{x_n\}_{n=1}^{N}$ denote a collection of $N$ audio clips, where each clip $x_n$ is itself a temporal signal composed of a sequence of acoustic frames. A clip's frame-level structure captures variations in background noise, speaker identity, accent, speaking rate, and recording quality. State-of-the-art ASR algorithms respond differently to these intra-clip variations \cite{li2014overviewrobust,cornell23_chime,bell2021adaptation}, and the best-performing algorithm for a given clip is generally not known in advance.

Let $\mathcal{A} = \{A_1, A_2, \ldots, A_K\}$ denote a finite set of $K$ pretrained ASR algorithms, where $A_j \in \mathcal{A}$ represents the $j$-th algorithm. For each clip $x_n$, the transcription produced by the $j$-th algorithm is $Y_{n,j} = A_j(x_n)$. Each transcription $Y_{n,j}$ is compared against the ground-truth reference $R_n$ using a transcription error measure. Using standard metrics such as the Word Error Rate (WER) \cite{li2014overviewrobust}, the transcription loss incurred by using algorithm $A_j$ on clip $x_n$ is $\ell(Y_{n,j}, R_n)$, where $\ell(\cdot,\cdot)$ denotes the chosen error metric.

Applying all $K$ algorithms to every incoming clip in order to pick the best transcription is computationally expensive. Instead, the goal is to select a single algorithm $j_n \in \{1,\dots,K\}$ for each clip such that the resulting transcription
$\hat{Y}_n = A_{j_n}(x_n)$ minimizes the cumulative transcription error across the dataset:
\begin{equation*}
\sum_{n=1}^{N} \ell(Y_{n,j_n}, R_n).
\end{equation*}

A key challenge is that each clip is internally non-stationary: speaker, channel, and background can shift across frames within the same clip, and these intra-clip dynamics influence which algorithm transcribes the clip most accurately. Selection mechanisms that compress a clip into a single pooled feature vector discard this temporal structure and lose information that is relevant to the choice of expert. We next introduce an algorithm that operates directly on frame-level encoder representations and uses a transformer-based router to model the temporal structure of each clip when picking an expert in an effective manner.


\section{HIT-ASR: Hierarchical Transformer Routing for Adaptive ASR Expert Selection}

\label{sec:workdone}

\begin{figure*}[t]
    \centering
    \includegraphics[width=1\linewidth]{figures/router_figure.pdf}
    \caption{Overview of the proposed adaptive ASR model selection
    framework. Each base ASR encoder produces a 2-D frame-level
    embedding tensor $\mathbf{H}_n^{(k)}\in\mathbb{R}^{T_n^{(k)}\times D_k}$
    with model-specific frame count $T_n^{(k)}$ and hidden dimension
    $D_k$; no temporal pooling is applied. Per-model linear
    projections, sinusoidal positional encodings, and learnable
    model-identity embeddings bring the heterogeneous sequences into
    a shared dimension. A shared Stage~1 transformer encoder
    summarizes each expert's frame sequence with a CLS token; an
    optional cross-attention bridge lets each summary attend to the
    frame-level outputs of the other experts; and a Stage~2
    transformer fuses the summaries through a global CLS token,
    which is mapped by a small MLP and softmax to the routing
    distribution $\mathbf{w}_n$. Only the argmax expert
    $A_{j_n}$ is decoded at inference.}
    \label{fig:architecture}
\end{figure*}

\subsection{Algorithm Overview}

HIT-ASR performs adaptive ASR model selection at the
clip level by exploiting the frame-level temporal structure of each
audio clip. The system operates in two phases: an offline training
phase and an online inference phase, as illustrated in
Fig.~\ref{fig:architecture}. The offline phase learns a
routing model that predicts, from the frame-level encoder
representations of a clip, which expert is most likely to transcribe
that clip with the lowest error. The online phase uses the trained
router to select an expert for each incoming clip and invokes only
that expert to produce the final transcription.

In the offline training phase, the system is provided with a
collection of labeled audio clips $\{(x_n,R_n)\}_{n=1}^{N}$. For each
clip $x_n$ and each expert $A_k\in\mathcal{A}$, the encoder of $A_k$
is run once on $x_n$ to obtain a frame-level hidden-state sequence
\[
\mathbf{H}_n^{(k)}=
\bigl(\mathbf{h}_{n,1}^{(k)},\dots,\mathbf{h}_{n,T_n^{(k)}}^{(k)}\bigr),
\qquad
\mathbf{h}_{n,t}^{(k)}\in\mathbb{R}^{D_k},
\]
where $T_n^{(k)}$ is the number of encoder frames produced by expert
$A_k$ for clip $x_n$ and $D_k$ is its hidden-state dimension. We use
$\mathbf{H}_n^{(k)}$ rather than a pooled summary so that intra-clip
acoustic dynamics are preserved. In parallel, the full ASR pipeline of
each expert is run on the clip to produce its candidate transcription
$Y_{n,k}=A_k(x_n)$ and the corresponding transcription error
$\mathrm{WER}(Y_{n,k},R_n)$. These errors serve as the supervision
signal that tells the router which expert was best for which clip.

The router takes the full set of frame-level hidden-state sequences
$\{\mathbf{H}_n^{(k)}\}_{k=1}^{K}$ as input and produces a probability
vector $\mathbf{w}_n=(w_{n,1},\dots,w_{n,K})$ over the experts. The
parameters of the router are optimized to align $\mathbf{w}_n$ with
the empirical per-clip error vector $(\mathrm{WER}(Y_{n,k},R_n))_{k=1}^{K}$,
using the training objective described in
Section~\ref{subsec:training_objective}.

In the online inference phase, the system receives a new audio clip
$x_n$. Each base encoder is run on $x_n$ to obtain its frame-level
hidden states $\mathbf{H}_n^{(k)}$, which are passed through the
trained router $g_\theta$ to produce $\mathbf{w}_n$. The expert with
the highest predicted probability,
\[
j_n=\arg\max_{k\in\{1,\dots,K\}} w_{n,k},
\]
is selected and only its decoder is executed on the clip to produce
the final transcription $\hat{Y}_n = A_{j_n}(x_n)$. Because only one
full ASR pipeline is decoded per clip, the proposed framework retains
the runtime cost of a single-model system on the decoding side, with
the additional cost of $K$ encoder forward passes and one router
forward pass.

The detailed components of the framework, including the frame-level
feature representation, the hierarchical routing architecture, and
the training procedure, are described in the following subsections.

\begin{algorithm}
\small
\caption{\small Online ASR Model Selection}
\label{alg:online_overview}

\textbf{Inputs:}
\begin{itemize}
    \item Audio clip $x_n$
    \item Set of base ASR encoders $\{E_1,\dots,E_K\}$ associated with the experts $\mathcal{A}=\{A_1,\dots,A_K\}$
    \item Trained routing function $g_{\theta}(\cdot)$
\end{itemize}

\textbf{Outputs:}
\begin{itemize}
    \item Final transcription $\widehat{Y}_n$ for clip $x_n$
\end{itemize}

\textbf{Inference Steps:}
\begin{algorithmic}[1]

\State For each expert, extract frame-level encoder representations:
\[
\mathbf{H}_n^{(k)} \leftarrow E_k(x_n), \qquad k=1,\dots,K
\]

\State Compute expert selection logits using the routing model:
\[
\mathbf{s}_n \leftarrow g_{\theta}\bigl(\{\mathbf{H}_n^{(k)}\}_{k=1}^{K}\bigr)
\]

\State Normalize logits using softmax to obtain routing probabilities $\mathbf{w}_n$

\State Select expert:
\[
j_n \leftarrow \arg\max_{k \in \{1,\dots,K\}} w_{n,k}
\]

\State Decode the selected expert to obtain the final transcription:
\[
\widehat{Y}_n \leftarrow A_{j_n}(x_n)
\]

\State Return $\widehat{Y}_n$

\end{algorithmic}
\end{algorithm}


\subsection{Frame-Level Feature Representation}
\label{sec:features}

For each clip $x_n$ and each expert $A_k$, we use the encoder of
$A_k$ as a feature extractor and keep its full frame-level
hidden-state sequence
\[
\mathbf{H}_n^{(k)}=
\bigl(\mathbf{h}_{n,1}^{(k)},\dots,\mathbf{h}_{n,T_n^{(k)}}^{(k)}\bigr),
\qquad
\mathbf{h}_{n,t}^{(k)}\in\mathbb{R}^{D_k}.
\]
Different encoders are based on different architectures and frame
rates, so $T_n^{(k)}$ and $D_k$ are model-specific. We do not pool
these frames into a single vector. Instead, the temporal sequence is
exposed directly to the router so that intra-clip dynamics, such as
shifts in noise level, speaker, or speaking rate, can influence the
routing decision.

This choice contrasts with prior model-selection approaches that
compress each clip into a single hand-crafted or pooled feature
vector. While such compact representations are inexpensive, they
discard the temporal structure that often distinguishes the regimes
under which different experts perform well. Keeping frame-level
sequences allows the router to learn, for example, that an expert is
preferable when a clip contains a stretch of overlapping speech or
when its spectral characteristics drift over time.

The frame-level features are obtained as a by-product of the encoder
forward pass that any expert would have to run anyway during
decoding. They therefore introduce no extra acoustic feature
extraction step beyond what the experts already perform.

\subsection{Hierarchical Transformer Router}
\label{subsec:transformer_router}

Given the per-expert frame-level sequences
$\{\mathbf{H}_n^{(k)}\}_{k=1}^{K}$ obtained from the feature pipeline
described in Section~\ref{sec:features}, the routing model has to map
these heterogeneous sequences (different $T_n^{(k)}$ and $D_k$ across
experts) to a single probability vector over the experts. We
implement the router as a hierarchical transformer with two stages.
The first stage produces a per-expert summary that captures the
intra-clip temporal structure as seen through that expert's encoder.
The second stage fuses these summaries across experts to produce the
final selection. An optional cross-attention bridge between the two
stages allows each per-expert summary to attend to the frame-level
representations of the other experts.

\subsubsection{Per-Model Projection and Positional Encoding}

Because the encoder dimension $D_k$ differs across experts, we first
project each frame to a common router dimension $d$ with a
model-specific linear transformation
\begin{equation}
\mathbf{p}_{n,t}^{(k)} = W_p^{(k)} \mathbf{h}_{n,t}^{(k)} + \mathbf{b}_p^{(k)},
\qquad W_p^{(k)} \in \mathbb{R}^{d\times D_k}.
\label{eq:projection}
\end{equation}
A standard sinusoidal positional encoding \cite{vaswani2017attention}
is added to each projected frame to preserve frame ordering within
the clip. To inform the router about which expert produced a given
frame, a learnable model-identity embedding $\mathbf{e}_k\in\mathbb{R}^{d}$
is added to every frame of expert $k$.

\subsubsection{Stage 1: Per-Expert Temporal Encoder}

For each expert $k$ we prepend a learnable summary token
$\mathbf{c}_k\in\mathbb{R}^{d}$ to its projected frame sequence,
forming
\[
\mathbf{X}_n^{(k)} =
\bigl(\mathbf{c}_k,\,\mathbf{p}_{n,1}^{(k)}+\mathbf{e}_k,\,\dots,\,
\mathbf{p}_{n,T_n^{(k)}}^{(k)}+\mathbf{e}_k\bigr).
\]
This sequence is processed by a Transformer encoder
\cite{vaswani2017attention}, denoted by $f_{\text{stage1}}$,
\begin{equation}
\mathbf{Y}_n^{(k)} = f_{\text{stage1}}\bigl(\mathbf{X}_n^{(k)}\bigr),
\label{eq:stage1}
\end{equation}
which performs masked multi-head self-attention and a position-wise
feed-forward transformation in each layer. Padding frames are excluded
through a key-padding mask. The output corresponding to the summary
token,
\[
\widetilde{\mathbf{c}}_n^{(k)} = \bigl[\mathbf{Y}_n^{(k)}\bigr]_{0},
\]
serves as a per-expert summary of clip $x_n$ that has been
temporally pooled by attention rather than by averaging.

By default the parameters of $f_{\text{stage1}}$ are shared across the
$K$ experts; this keeps the router compact and lets the model-identity
embedding $\mathbf{e}_k$ carry the expert-specific information. Using
separate Stage~1 parameters per expert is also supported, and is
evaluated as part of the architecture ablation in
Section~\ref{sec:experiments}.

\subsubsection{Cross-Attention Bridge (Optional)}

Stage~1 produces the per-expert summary $\widetilde{\mathbf{c}}_n^{(k)}$
using only the frames of expert $k$. To let each summary be informed
by fine-grained frame-level evidence from the \emph{other} experts,
we optionally insert a cross-attention bridge. For each expert $k$,
let
\[
\mathbf{F}_n^{(\neg k)} = \bigl[\mathbf{Y}_n^{(j)}\bigr]_{1{:}T_n^{(j)}}\,\big|\,j\neq k
\]
denote the concatenated frame-level outputs of all other experts (with
their CLS positions removed). The bridge updates each summary by
\begin{equation}
\widetilde{\mathbf{c}}_n^{(k)} \leftarrow
\mathrm{LN}\Bigl(
\widetilde{\mathbf{c}}_n^{(k)}
+ \mathrm{MHA}\bigl(
\widetilde{\mathbf{c}}_n^{(k)},\,\mathbf{F}_n^{(\neg k)},\,\mathbf{F}_n^{(\neg k)}
\bigr)
\Bigr),
\label{eq:bridge}
\end{equation}
where $\mathrm{MHA}$ is multi-head attention with the summary as
query and the other-experts' frames as keys and values, and
$\mathrm{LN}$ denotes layer normalization. Because each query is a
single token, this step is inexpensive relative to Stage~1.

\subsubsection{Stage 2: Cross-Expert Fusion}

The (possibly bridged) per-expert summaries
$\bigl\{\widetilde{\mathbf{c}}_n^{(k)}\bigr\}_{k=1}^{K}$ are
concatenated together with a learnable global summary token
$\mathbf{c}^{\star}\in\mathbb{R}^{d}$ to form the Stage~2 input
sequence
\[
\mathbf{Z}_n =
\bigl(\mathbf{c}^{\star},\,
\widetilde{\mathbf{c}}_n^{(1)},\dots,\widetilde{\mathbf{c}}_n^{(K)}\bigr).
\]
A second Transformer encoder $f_{\text{stage2}}$ processes
$\mathbf{Z}_n$ to mix information across experts. The updated global
token, $\widetilde{\mathbf{c}}^{\star}_n$, is fed to a small MLP
classifier to produce the routing logits
\begin{equation}
\mathbf{s}_n =
W_2\,\mathrm{GELU}\bigl(W_1\widetilde{\mathbf{c}}^{\star}_n
+\mathbf{b}_1\bigr) + \mathbf{b}_2 \in \mathbb{R}^{K}.
\label{eq:logits}
\end{equation}
The corresponding routing probabilities are obtained by a softmax
\begin{equation}
w_{n,k} =
\frac{\exp(s_{n,k})}{\sum_{m=1}^{K}\exp(s_{n,m})},
\qquad k=1,\dots,K.
\label{eq:router_softmax}
\end{equation}
The selected expert is the one with the highest predicted
probability, as summarized in Algorithm~\ref{alg:online_overview}.

\subsection{Training Objective}
\label{subsec:training_objective}

The router parameters $\theta$ are learned offline from a labeled
dataset
\[
\mathcal{D}=\{(x_n,R_n)\}_{n=1}^{N}
\]
of audio clips and their reference transcriptions. Before training,
each expert $A_k$ is run on every clip to obtain the candidate
transcription $Y_{n,k}=A_k(x_n)$ and the per-clip Word Error Rate
\[
\mathrm{WER}_{n,k} \;=\; \mathrm{WER}(Y_{n,k},R_n)
\;=\; \frac{S+I+D}{|R_n|},
\]
where $S$, $I$, and $D$ are the substitution, insertion, and deletion
counts from a Levenshtein alignment \cite{li2014overviewrobust}, and
$|R_n|$ is the number of words in the reference. Encoder hidden
states $\{\mathbf{H}_n^{(k)}\}$ are cached at the same time. During
router training, both the WER scores and the encoder features are
treated as fixed inputs.

We combine three loss terms. The first is a primary
\emph{weighted-WER} loss that directly minimizes the expected
transcription error under the router's distribution:
\begin{equation}
\mathcal{L}_{\text{wer}}(\theta) =
\frac{1}{N}\sum_{n=1}^{N}\sum_{k=1}^{K}
w_{n,k}\,\mathrm{WER}_{n,k}.
\label{eq:loss_wer}
\end{equation}
This loss directly aligns the router with task-level performance, but
the gradient signal that any single sample provides is weak when all
$\mathrm{WER}_{n,k}$ values are similar in magnitude.

To strengthen the supervision early in training, we add an auxiliary
\emph{hard cross-entropy} term against the per-clip oracle expert,
$k_n^\star=\arg\min_{k}\mathrm{WER}_{n,k}$, with label smoothing
$\varepsilon\in[0,1)$:
\begin{equation}
\mathcal{L}_{\text{hard}}(\theta) =
-\frac{1}{N}\sum_{n=1}^{N}\sum_{k=1}^{K}
\tilde{q}_{n,k}\,\log w_{n,k},
\label{eq:loss_hard}
\end{equation}
where $\tilde{q}_{n,k}=(1-\varepsilon)\mathbb{1}[k=k_n^\star]+\varepsilon/K$.
This term provides a stronger gradient but ignores how much better the
oracle expert is compared to the others.

To preserve the WER-magnitude information while keeping a strong
gradient, we additionally use a \emph{soft cross-entropy} term
against a temperature-controlled target distribution derived from the
WER vector,
\begin{equation}
q_{n,k}^{\mathrm{soft}} =
\frac{\exp(-\mathrm{WER}_{n,k}/\tau)}
{\sum_{m=1}^{K}\exp(-\mathrm{WER}_{n,m}/\tau)},
\label{eq:soft_target}
\end{equation}
\begin{equation}
\mathcal{L}_{\text{soft}}(\theta) =
-\frac{1}{N}\sum_{n=1}^{N}\sum_{k=1}^{K}
q_{n,k}^{\mathrm{soft}}\,\log w_{n,k}.
\label{eq:loss_soft}
\end{equation}
Small $\tau$ recovers a near-hard target sharply concentrated on the
oracle expert; larger $\tau$ produces a smoother target that reflects
the relative WER gaps between experts.

The overall training objective is a weighted sum
\begin{equation}
\min_{\theta}\;
\mathcal{L}(\theta) =
\lambda_{\text{wer}}\,\mathcal{L}_{\text{wer}}(\theta)
+\lambda_{\text{hard}}\,\mathcal{L}_{\text{hard}}(\theta)
+\lambda_{\text{soft}}\,\mathcal{L}_{\text{soft}}(\theta),
\label{eq:total_objective}
\end{equation}
with non-negative coefficients
$(\lambda_{\text{wer}},\lambda_{\text{hard}},\lambda_{\text{soft}})$
that determine which combination of the three signals is used. The
choice of these coefficients is studied in the loss ablation in
Section~\ref{sec:experiments}.

\vspace{0.1cm}
\noindent\textbf{Remark.}\,\textit{Although WER is used as the
primary error measure due to its widespread use in speech recognition
research~\cite{yu2016asrdeep,li2014overviewrobust}, any non-negative
per-clip error metric, such as Character Error Rate or a
domain-specific score, can be substituted into
\eqref{eq:loss_wer}--\eqref{eq:loss_soft} without modifying the
router architecture.}
\vspace{0.1cm}

The complete offline training procedure is summarized in
Algorithm~\ref{alg:offline_training}.

\begin{algorithm}
\small
\caption{Offline Training of the Hierarchical Router}
\label{alg:offline_training}

\textbf{Inputs}
\begin{itemize}
\item $\{(x_n,R_n)\}_{n=1}^{N}$: Labeled dataset of audio clips and references
\item $\mathcal{A}=\{A_1,\dots,A_K\}$ and their encoders $\{E_1,\dots,E_K\}$
\item Loss coefficients $(\lambda_{\text{wer}},\lambda_{\text{hard}},\lambda_{\text{soft}})$, label smoothing $\varepsilon$, temperature $\tau$
\end{itemize}

\textbf{Output} \\
Trained router $g_{\theta}$

\textbf{Training Steps}
\begin{algorithmic}[1]
\State Split dataset into training and validation sets:
$\mathcal{D}_{\text{train}}, \mathcal{D}_{\text{val}}$

\For{$n \gets 1$ \textbf{to} $N$}
    \For{$k \gets 1$ \textbf{to} $K$}
        \State Cache encoder features:
        $\mathbf{H}_n^{(k)} \gets E_k(x_n)$
        \State Decode and score:
        $Y_{n,k} \gets A_k(x_n);\;
        \mathrm{WER}_{n,k}\gets\mathrm{WER}(Y_{n,k},R_n)$
    \EndFor
\EndFor

\For{each training step}
    \State Sample a minibatch $\mathcal{B}\subset\mathcal{D}_{\text{train}}$
    \State Compute routing probabilities
    $\mathbf{w}_n$ from $\{\mathbf{H}_n^{(k)}\}$ using
    \eqref{eq:projection}--\eqref{eq:router_softmax}
    \State Compute soft target $q^{\mathrm{soft}}_n$ via \eqref{eq:soft_target}
    \State Compute composite loss $\mathcal{L}(\theta)$ on $\mathcal{B}$ via \eqref{eq:total_objective}
    \State Update $\theta$ with a gradient-based optimizer (e.g.\ AdamW)
\EndFor

\State Perform model selection and early stopping using
$\mathcal{L}_{\text{wer}}$ on $\mathcal{D}_{\text{val}}$
\end{algorithmic}
\end{algorithm}

\section{Experiments}
\label{sec:experiments}

This section evaluates the proposed router along three axes. First, a
synthetic experiment isolates the contribution of frame-level
temporal modeling under a controlled regime-switching setting
(Section~\ref{subsec:synthetic}). Second, real-world routing on AMI
and VoxPopuli compares the router against the individual base
models and a set of selection and combination baselines
(Section~\ref{subsec:real}). Third, ablations on the loss and the
router architecture identify which design choices drive the observed
gains (Section~\ref{subsec:ablation}). Section~\ref{subsec:discussion}
discusses what the experiments support and where the framework breaks
down.

% Throughout the section we report Word Error Rate ($\mathrm{WER}$) as
% the primary metric and, where applicable, the routing-specific
% quantities introduced in Section~\ref{sec:workdone}: \emph{selection
% accuracy} (the fraction of clips for which the router picks the
% oracle expert) and the \emph{oracle gap}
% $\mathrm{WER}_{\text{router}}-\mathrm{WER}_{\text{oracle}}$, where
% $\mathrm{WER}_{\text{oracle}}$ is the per-clip minimum WER across
% experts. The oracle gap measures how far the router is from the
% hindsight-optimal selection on the same clips.

\subsection{Synthetic Regime-Switch Experiment}
\label{subsec:synthetic}

\subsubsection{Motivation}

The central design choice in this work is to keep the per-expert
encoder outputs at frame-level resolution rather than collapsing them
into a single pooled vector. A useful sanity check is therefore a
controlled setting in which the optimal expert for a clip is
\emph{determined by the temporal structure of the clip itself} rather
than by any pooled summary of it. If the proposed router exploits
temporal information, it should outperform a router that has access
to the same per-expert features but only after temporal pooling. If
both routers perform similarly, the temporal-modeling claim is not
supported.

\subsubsection{Synthetic Data Generation}

We construct a synthetic dataset of $N=10{,}000$ clips, each with a
fixed length of $T=128$ frames. Each clip $x_n$ is divided into two
halves separated by a regime-switch at frame $\rho_n\sim\mathcal{U}\{T/4,3T/4\}$.
The first half is drawn from one acoustic regime $r_n^{(1)}$ and the
second half from another regime $r_n^{(2)}\neq r_n^{(1)}$; the two
regimes are sampled uniformly from a finite set of $R=4$ regimes.
For each regime we fix a center vector $\boldsymbol{\mu}_r$, and each
frame is generated as
\[
\mathbf{h}_{n,t}^{(k)} =
W_k\bigl(\boldsymbol{\mu}_{r_n^{(t)}}+\boldsymbol{\eta}_{n,t}\bigr),
\qquad
\boldsymbol{\eta}_{n,t}\sim\mathcal{N}(\mathbf{0},\sigma^2 I),
\]
where $W_k$ is an expert-specific projection that simulates the
hidden-state space of expert $A_k$ and $r_n^{(t)}=r_n^{(1)}$ for
$t\le\rho_n$ and $r_n^{(2)}$ otherwise. We use $K=3$ experts.

We assign per-expert WER scores via a fixed table
$\mathrm{WER}_{r_1,r_2,k}$ that depends only on the \emph{ordered
pair} of regimes within the clip. The table is constructed so that
no expert is uniformly best across regime pairs, the best expert
depends on \emph{both} halves (so a router that sees only the
average regime cannot recover it), and a small Gaussian noise is
added to each $\mathrm{WER}_{n,k}$ to make the optimization
non-trivial.

\subsubsection{Models Compared}

We compare four selection strategies on the synthetic data:
\begin{itemize}
    \item Random -- expert chosen uniformly at random per clip.
    \item MLP-pool -- a router that mean-pools each per-expert frame
    sequence to a single vector and feeds the concatenation through
    a 2-layer MLP. Same parameter budget as the proposed router,
    same training loss.
    \item Proposed -- the hierarchical transformer router of
    Section~\ref{subsec:transformer_router}.
    \item Oracle -- selects the per-clip
    $\arg\min_k\mathrm{WER}_{n,k}$.
\end{itemize}
Both learned routers are trained with the same loss
\eqref{eq:total_objective} and the same optimizer, on a $\,80/10/10$
train/validation/test split.

\subsubsection{Results}

Table~\ref{tab:synthetic} reports the average WER, selection
accuracy, and oracle gap on the test split.
The expected pattern is that the MLP-pool router improves over Random
but plateaus well above Oracle, while the proposed router perfectly 
identifies the regime-switch position from the frame sequence directly, 
achieving an oracle gap of zero. We additionally vary the regime-switch position
$\rho_n$ and the number of regimes $R$ in
Fig.~\ref{fig:synthetic_sweep} to show how the temporal-modeling
advantage scales with the difficulty of the temporal pattern.

\begin{table}[t]
\centering
\caption{Word Error Rate (WER) and selection accuracy results on the synthetic regime-switch dataset. The proposed HIT-ASR perfectly identifies the regime boundaries, matching the theoretical Oracle performance.}
\label{tab:synthetic}
\small
\begin{tabular}{l S[table-format=2.1] S[table-format=3.1]}
\toprule
\textbf{Method} & {\textbf{WER (\%, $\downarrow$)}} & {\textbf{Sel.\ acc.\ (\%, $\uparrow$)}} \\
\midrule
Random                        & 37.9 &  33.2 \\
Weighted random               & 37.5 &  34.3 \\
MLP-pool                      & 28.9 &  58.1 \\
HIT-ASR                       & 15.3 & 100.0 \\
\midrule
Oracle                        & 15.3 & 100.0 \\
\bottomrule
\end{tabular}
\end{table}

\subsection{Real-World Results}
\label{subsec:real}

\subsubsection{Datasets}

We evaluate on two standard long-form English speech corpora:
\begin{itemize}
    \item AMI~\cite{carletta2005ami}: meeting recordings with multiple
    speakers, overlapping speech, and varying microphone conditions.
    We segment the audio into utterance-level clips using the provided annotations.
    \item VoxPopuli~\cite{wang2021voxpopuli}: parliamentary speeches
    with diverse speakers, accents, and recording channels. We use
    the English subset and follow the published splits.
\end{itemize}
Both datasets are pre-processed to obtain per-clip references, and
each base ASR model is run once over the entire data to cache its
encoder hidden states and its candidate transcription. The router is
then trained on the cached features.

\subsubsection{Base ASR Models}

We use $K=3$ pretrained ASR models with different backbones:
\begin{itemize}
    \item HuBERT-Large~\cite{hsu2021hubert}, fine-tuned on
    LibriSpeech;
    % (\texttt{facebook/hubert-large-ls960-ft}); 
    encoder dimension $D=1024$.
    \item Whisper-base~\cite{radford2022whisper};
    % (\texttt{openai/whisper-base}); 
    encoder dimension $D=512$.
    \item Wav2Vec2-Large-Robust~\cite{baevski2020wav2vec}, fine-tuned
    on Switchboard;
    % (\texttt{facebook/wav2vec2-large-robust-ft-swbd-300h}); 
    encoder
    dimension $D=1024$.
\end{itemize}
These three models are trained on different data and have different
strengths; AMI and VoxPopuli are out-of-distribution for at least one
of them, which is the regime where adaptive selection is expected to
help.

\subsubsection{Baselines}

We compare the proposed router against:
\begin{itemize}
    \item Single base model (HuBERT, Whisper, Wav2Vec2 separately): no
    selection, runs the same model on every clip.
    \item Random: picks an expert uniformly at random per clip.
    \item Weighted random: samples an expert from the prior
    proportional to inverse training-set WER.
    \item ROVER~\cite{fiscus1997rover}: hypothesis-level voting that
    runs all $K$ experts and combines their outputs.
    \item MLP-pool selector: same as in
    Section~\ref{subsec:synthetic}; mean-pools each expert's
    frame-level features and feeds the concatenation to an MLP
    classifier, trained with the same loss as the proposed router.
    \item Oracle: per-clip $\arg\min_k\mathrm{WER}_{n,k}$, an upper
    bound on any selector that picks one of the $K$ experts.
\end{itemize}
The single-model and ROVER baselines bound what can be achieved
without learning to route, and the MLP-pool baseline isolates the
contribution of the frame-level transformer relative to a purely
clip-pooled learner.

\subsubsection{Training Setup}

The router is trained with AdamW (learning rate $10^{-4}$, weight
decay $10^{-2}$), linear warmup of $200$ steps followed by cosine
annealing, batch size $8$, gradient clipping at $1.0$, and mixed
precision. The default hyperparameters are $d=256$, $n_{\text{heads}}=4$,
$\ell_{\text{stage1}}=2$, $\ell_{\text{stage2}}=1$,
$d_{\text{ffn}}=512$, dropout $0.15$. Frame sequences longer than
$T_{\max}=2000$ are truncated. Cross-attention bridge is enabled and
Stage~1 weights are shared across experts unless stated otherwise.
For the main results we use the best-performing loss configuration
identified in Section~\ref{subsec:ablation}, namely
$(\lambda_{\text{wer}},\lambda_{\text{hard}},\lambda_{\text{soft}})=(1,0,0.5)$
with soft-CE temperature $\tau=0.1$. Early stopping uses
$\mathrm{WER}_{\text{router}}$ on validation with patience $10$. To ensure rigorous evaluation and prevent data leakage, all reported metrics for AMI and VoxPopuli represent the held-out test performance averaged across a 5-fold cross-validation setup. The router is trained exclusively on the respective training splits in each fold.

\subsubsection{Results}

Table~\ref{tab:main_results} reports test WER on AMI and VoxPopuli. Transitioning from the synthetic experiment to real-world datasets reveals an expected increase in the oracle gap (e.g., from zero to $\sim$10\% absolute WER on AMI). While the synthetic regimes are perfectly demarcated by clean Gaussian shifts, real-world acoustics involve complex, overlapping variations and potential feature extraction bottlenecks that make ideal routing more challenging. 

Notably, Whisper-base exhibits an abnormally high WER (121.40\%) on VoxPopuli due to severe hallucination loops on certain utterances. However, it remains the most accurate expert for roughly 40\% of the individual clips. HIT-ASR effectively navigates this volatility; by relying on the frame-level evidence, the router learns to avoid Whisper on clips prone to hallucination while still exploiting it when it is optimal.

\begin{table}[t]
\centering
\caption{Word Error Rate (\%, $\downarrow$) on the real-world AMI and VoxPopuli datasets. The Oracle represents the theoretical lower bound if the best expert was selected for every clip in hindsight.}
\label{tab:main_results}
\small
\begin{tabular}{l S[table-format=3.2] S[table-format=3.2]}
\toprule
\textbf{Model} & {\textbf{AMI (\%)}} & {\textbf{VoxPopuli (\%)}} \\
\midrule
HuBERT-Large & 64.92 & 23.16 \\
Whisper-base & 57.68 & 121.40 \\
Wav2Vec2-Large-Robust & 38.52 & 29.81 \\
\midrule
Random & 54.44 & 56.57 \\
Weighted random & 49.22 & 26.28 \\
ROVER & 43.10 & 24.52 \\
MLP-pool (hard CE) & 38.31 & 29.81 \\
HIT-ASR & \textbf{34.23} & \textbf{23.04} \\
\midrule
Oracle & 24.13 & 16.47 \\
\bottomrule
\end{tabular}
\end{table}

\subsection{Ablation Study}
\label{subsec:ablation}

\subsubsection{Loss Ablation}

The training objective \eqref{eq:total_objective} combines three
signals -- weighted-WER, hard cross-entropy on the oracle expert, and
a temperature-controlled soft cross-entropy. Table~\ref{tab:loss_ablation}
isolates the effect of each by sweeping the coefficients
$(\lambda_{\text{wer}},\lambda_{\text{hard}},\lambda_{\text{soft}})$
and $\tau$ on AMI, keeping all other hyperparameters fixed.

\begin{table}[t]
\centering
\caption{Impact of the loss function components on the AMI test set. The soft cross-entropy (CE) term provides critical Word Error Rate (WER) magnitude information to the router. The hyperparameter $\tau$ denotes the soft-CE temperature.}
\label{tab:loss_ablation}
\small
\begin{tabular}{l S[table-format=1.1] S[table-format=1.1] S[table-format=1.1] c S[table-format=2.1]}
\toprule
\textbf{Loss config.} & {$\lambda_{\text{wer}}$} & {$\lambda_{\text{hard}}$} & {$\lambda_{\text{soft}}$} & {$\tau$} & {\textbf{WER (\%, $\downarrow$)}}\\
\midrule
WER only                 & 1.0 & 0.0 & 0.0 & {--}  & 38.2 \\
Hard CE only             & 0.0 & 1.0 & 0.0 & {--}  & 37.1 \\
Soft CE only             & 0.0 & 0.0 & 1.0 & 0.1 & 34.3 \\
Soft CE only             & 0.0 & 0.0 & 1.0 & 1.5 & 33.5 \\
WER $+$ Hard CE          & 1.0 & 0.3 & 0.0 & {--}  & 34.7 \\
WER $+$ Soft CE          & 1.0 & 0.0 & 0.5 & 0.1 & 34.2 \\
WER $+$ Soft CE          & 1.0 & 0.0 & 0.5 & 1.5 & 35.6 \\
Heavy Soft CE            & 1.0 & 0.0 & 1.0 & 0.5 & 33.5 \\
Full Recipe              & 1.0 & 0.3 & 0.5 & 0.5 & 33.9 \\
Full Recipe              & 1.0 & 0.3 & 0.5 & 1.5 & \textbf{33.4} \\
\midrule
Oracle                   & {--}  & {--}  & {--}  & {--}  & 23.7 \\
\bottomrule
\end{tabular}
\end{table}

Two patterns are worth noting. First, optimizing strictly to match the oracle expert does not align perfectly with minimizing the final WER: \emph{Hard CE only} and \emph{WER $+$ Hard CE} penalize any deviation from the oracle, but perform poorly in terms of overall WER. This occurs because picking the oracle expert is rewarded equally regardless of how much better the oracle was compared to the alternatives. Second, variants incorporating \emph{Soft CE} achieve the lowest WER: the soft target carries WER-magnitude information into the gradient, which steers the router toward decisions whose mistakes are cheap, even if it occasionally misses the exact oracle. The full recipe and its variants balance these trade-offs, establishing a clean story that optimizing for the oracle is not identical to optimizing for WER. This is the configuration we report as the default in Section~\ref{subsec:real}.

\subsubsection{Architecture Ablation}

The router architecture has three primary design switches: the
cross-attention bridge, the sharing of Stage~1 weights across
experts, and the depth of each stage. Table~\ref{tab:arch_ablation}
toggles each on AMI.

Notably, removing the cross-attention bridge and parameter sharing 
simultaneously achieves the best performance ($33.6\%$ WER). However, 
this is expected because removing parameter sharing substantially 
increases the model capacity (from $4.6$M to $6.7$M parameters), 
compensating for the lack of a bridge. The default configuration 
retains parameter sharing as a strong regularizer to maintain a 
lightweight footprint, and under this shared regime, the 
cross-attention bridge is critical, improving the WER from $34.9\%$ 
to $34.4\%$ without a massive parameter increase.

\begin{table}[t]
\centering
\caption{Ablation of the HIT-ASR router architecture on the AMI test set. The default configuration uses shared Stage 1 weights and a cross-attention (CA) bridge to balance capacity and regularization.}
\label{tab:arch_ablation}
\small
\begin{tabular}{p{4.1cm} S[table-format=2.1] S[table-format=2.1]}
\toprule
\textbf{Variant} & {\textbf{WER (\%, $\downarrow$)}} & {\textbf{Params (M)}} \\
\midrule
\textit{Reference} & & \\
Default (shared S1 + bridge) & 34.4 & 4.6 \\
\midrule
\textit{Connectivity / sharing} & & \\
$-$ CA bridge & 34.9 & 3.8 \\
$-$ shared Stage~1 (per expert) & 34.2 & 7.5 \\
$-$ CA bridge \& shared Stage~1 & 33.6 & 6.7 \\
\midrule
\textit{Depth} & & \\
$\ell_{\text{stage1}}=1$ & 34.3 & 3.0 \\
$\ell_{\text{stage1}}=4$ & 34.5 & 7.8 \\
$\ell_{\text{stage2}}=2$ & 34.1 & 5.2 \\
\midrule
\textit{Width / heads / regularization} & & \\
$d=128$ & 35.0 & 1.4 \\
$d=512$ & 34.4 & 14.9 \\
$h=2$ & 33.9 & 4.6 \\
$h=8$ & 34.4 & 4.6 \\
$\text{dropout}=0.2$ & 34.9 & 4.6 \\
\bottomrule
\end{tabular}
\end{table}


\subsection{Discussion}
\label{subsec:discussion}

The synthetic experiment isolates the temporal-modeling claim and
shows that frame-level routing outperforms a same-budget pooled
router when the regime-switch is the only source of expert-selection
information. The real-world experiments translate this into
end-to-end WER reductions on AMI and VoxPopuli relative to the best
single base model and the ROVER and MLP-pool baselines, while only
running one full ASR pipeline per clip at inference time.

The loss ablation surfaces a non-obvious point: optimizing the
selection accuracy and optimizing the resulting WER are not the
same objective. The soft cross-entropy term, by carrying the WER
gaps directly into the gradient, drives the router toward decisions
whose worst case is bounded, even when it sometimes \emph{misses}
the oracle expert in cases where the gap is small. This matches the
intuition that what matters operationally is how bad the chosen
expert is, not how often it equals the oracle.

The proposed framework has two limitations worth stating. First, the
router still requires the encoder of every base model to run on each
clip, even though only one full ASR pipeline is decoded; the savings
come from the decoder side. For models with comparable encoder and
decoder costs, the inference saving is correspondingly smaller.
Second, the router operates per clip and does not share information
across clips of the same recording. Long-form scenarios where speaker
and channel persist across clips are likely to benefit from
across-clip context, which we leave as a natural extension.

\section{Conclusion}
\label{sec:conclusion}

In this work, we presented HIT-ASR, a feature-adaptive model selection framework designed to route diverse audio clips to the most suitable pretrained ASR expert. We demonstrated that preserving the temporal structure of an audio clip through frame-level encoder representations is crucial for making accurate routing decisions, particularly in long-form and heterogeneous audio streams where acoustic conditions fluctuate unpredictably. Our hierarchical transformer router effectively models this intra-clip temporal evidence, and its integration with a soft cross-entropy loss ensures that the system directly optimizes for transcription quality (WER) rather than mere categorical selection accuracy. Empirical results on both synthetic regime-switch scenarios and complex real-world datasets (AMI and VoxPopuli) confirm that HIT-ASR yields significant WER reductions over static models, pooling-based routers, and hypothesis-level combination strategies like ROVER. Moreover, by decoding only with the selected expert, the framework maintains an inference cost comparable to a single ASR model. Future work will explore extending this temporal modeling paradigm to capture cross-clip contextual dependencies for extended conversational sessions and integrating more heterogeneous pools of base experts.

% \clearpage

\bibliographystyle{ieeetr}
\bibliography{reference}





\end{document}


% %In order to address the online model selection problem defined in Section II, we propose a novel randomized switching-based meta-learning algorithm for real-time STT transcription. Unlike static ensemble methods, which aggregate predictions across all models, or deterministic selectors, which may overfit to specific conditions, our method dynamically selects a single transcription model at each time step while adaptively tracking nonstationary changes in the audio environment. The algorithm is inspired by the universal randomized switching framework and is designed to achieve sublinear regret with respect to the best piecewise-constant switching sequence in hindsight.


% \subsection*{Randomized Switching-Based Selection}

% To achieve this, we maintain a set of probability weights $\{w_{t,m}\}_{m=1}^K$ that represent the likelihood of selecting each model at time $t$. These probabilities are updated adaptively based on past losses and compactly encode all possible switching paths without requiring exponential enumeration.

% At each step $t$:
% \begin{enumerate}
%     \item \textbf{Sampling:} We sample a model index $m_n$ from a categorical distribution defined by the weight vector $\mathbf{w}_t = (w_{t,1},\ldots,w_{n,k})$:
%     \[
%     m_n \sim \mathrm{Cat}(\mathbf{w}_t), \quad \Pr[m_n = m] = w_{t,m}.
%     \]
%     The selected model produces the transcription $T_{n,m_n}$.

%     \item \textbf{Loss Observation:} The algorithm observes $\ell(T_{n,m_n}, R_n)$ for the chosen model. In a bandit setting, importance-weighted estimators can be used to obtain unbiased estimates for unobserved losses.

%     \item \textbf{Auxiliary State Update:} For each possible last-switch time $\ell \leq t$, we maintain auxiliary variables $Z_t(\ell,m)$ that accumulate exponentiated loss contributions. These variables are updated recursively according to:
%     \[
%     Z_t(\ell,m) = \beta Z_{t-1}(\ell,m) \cdot \exp(-\eta \ell(T_{n,m},R_n)) + \frac{1-\beta}{K} \sum_{j=1}^K Z_{t-1}(\ell,j),
%     \]
%     where $\eta > 0$ is a learning rate and $\beta \in (0,1)$ controls the switching penalty.

%     \item \textbf{Weight Update:} The normalized weights are then computed as
%     \[
%     w_{t+1,m} = \frac{\sum_{\ell=1}^{t} Z_t(\ell,m)}{\sum_{j=1}^K \sum_{\ell=1}^{t} Z_t(\ell,j)}.
%     \]
% \end{enumerate}

% \subsection*{Algorithmic Properties}

% The above procedure ensures that the selection strategy:
% \begin{itemize}
%     \item Tracks the best switching sequence in hindsight with expected regret bounded by $O(\sqrt{T \log T})$.
%     \item Operates with $O(K)$ update complexity per time step, making it efficient for real-time STT applications.
%     \item Naturally adapts to nonstationary environments where speaker conditions or noise profiles change abruptly or gradually.
% \end{itemize}

% \subsection*{Integration with STT Model Selection}

% In practice, each STT algorithm $M_m$ can be a large pretrained STT model (e.g., Whisper, wav2vec2, RNN-T, or domain-adapted fine-tunes). The randomized switching algorithm provides a principled way to select among them on a per-utterance basis. The stochastic nature of selection ensures continuous exploration of all candidate models, while the adaptive weight updates concentrate probability mass on those models that demonstrate superior recent performance. This framework thus provides both theoretical guarantees and practical efficiency for real-time deployment in adaptive STT systems.
